\chapter{ANALISIS MASALAH DAN RANCANGAN SOLUSI}
\label{chap:analisis-solusi}
\begin{petunjuk}
Tujuan utama penulisan bab ini adalah untuk menguraikan masalah dan gambaran solusi penyelesaian masalah tugas akhir spesifik individu anggota tim capstone. Judul bab 3 dapat disesuaikan dengan konteks persoalan, misalnya: Analisis Masalah dan Rancangan Solusi Pembangunan Sistem ABC. Bab ini mencakup antara lain subbab-subbab:
\end{petunjuk}

\section{Analisis Masalah}
\begin{petunjuk}
Deskripsi dan analisis dikaitkan dengan Rumusan Masalah tugas akhir, misalnya menjelaskan secara detail latar belakang dan masalah atau tantangan yang menjadi dasar perlunya solusi tugas akhir, menunjukkan \textit{gap}/celah antara kondisi saat ini dengan kondisi yang diharapkan, serta dapat juga dijelaskan kaitan antara sistem/aplikasi yang dikembangkan dengan sistem/aplikasi lain yang terkait.
\end{petunjuk}

\section{Gambaran Umum Solusi}
\begin{petunjuk}
Analisis solusi dapat terdiri dari pilihan alternatif solusi masalah tugas akhir untuk setiap permasalahan berdasarkan hasil studi literatur atau survei, pemilihan solusi beserta justifikasinya.

Deskripsi umum gambaran awal solusi tugas akhir yang dipilih, antara lain dapat mencakup:
\begin{itemize}
  \item Gambaran awal modul/subsistem/komponen solusi yang akan dikembangkan untuk menyelesaikan masalah sebagai bagian dari solusi tugas akhir, berikut penjelasannya.
  \item Gambaran awal alur umum algoritma atau langkah-langkah pengembangan solusi tugas akhir dan penjelasannya.
  \item Penggunaan kakas secara umum yang diperlukan.
\end{itemize}
Untuk memperjelas gambaran solusi, dianjurkan menggunakan gambar diagram sebagai pendukung penjelasan bagian ini.
\end{petunjuk}
